\section{Shannon's information theory}

Shannon's information theory provides a mathematical framework for understanding communication and quantifying information.

Communication occurs between two endpoints:
\begin{itemize}
    \item The sender comprises an information source and an encoder.
    \item The receiver consists of an information destination and a decoder.
\end{itemize}
Information is transmitted through a channel in the form of a sequence of symbols from a finite alphabet.
\begin{figure}[H]
    \centering
    \includegraphics[width=1\linewidth]{images/sha.png}
    \caption{Shannon's communication structure}
\end{figure}
The receiver exclusively receives information through the channel. 
Until the symbol arrives, there remains uncertainty about what the next symbol will be. 
Consequently, we represent the sender as a random variable. 
Hence, obtaining information is akin to determining an outcome of a random variable $\mathcal{X}$, and the quantity of information relies on the distribution of $\mathcal{X}$. 
Encoding involves mapping each outcome to a finite sequence of symbols: more symbols are necessary when transmitting more information.

\subsection{Entropy}
We require a non-negative measure of uncertainty. 
The combination of uncertainties should correspond to adding entropies.
\begin{definition}[\textit{Entropy}]
    Let $\mathcal{X}$ be a discrete random variable with $n$ outcomes in $\left\{ x_0,\dots,x_{n-1} \right\}$, where $\Pr(\mathcal{X}=x_i)=p_i$ for all $0 \leq 1 \leq n$. 
    The entropy of $\mathcal{X}$ is given by:
    \[H(\mathcal{X})=\sum_{i=0}^{n-1}-\Pr_i\log_b\left(\Pr_i\right)\]
\end{definition}
The unit of measurement for entropy is contingent on the base $b$ of the logarithm, where the typical case for $b = 2$ is bits.
\begin{example}
    The random variable $\mathcal{X}$ represents a sequence of 6 uniform random letters (with $6^{26}$ combinations). 
    In this case, the entropy is calculated as:
    \[H(\mathcal{X})=\sum_{i=0}^{6^{26}-1}-\dfrac{1}{6^{26}}\log_b\left(\dfrac{1}{6^{26}}\right)\approx 28.2b\]

    On the other hand, if $\mathcal{X}$ represents a uniform selection from six-letter English words (with $6^{6}$ combinations), the entropy is computed as:
    \[H(\mathcal{X})=\sum_{i=0}^{6^{6}-1}-\dfrac{1}{6^{6}}\log_b\left(\dfrac{1}{6^{6}}\right)\approx 12.6b\]
\end{example}

\begin{theorem}
    It is possible to encode the outcomes $n$ of independent and identically distributed random variables, each with entropy $H(\mathcal{X})$, using at least $nH(\mathcal{X})$ bits per outcome. 
    Encoding with fewer than $nH(\mathcal{X})$ bits will result in loss of information.
\end{theorem}
Consequently, achieving arbitrary compression of bit-strings without loss is unattainable, necessitating cryptographic hashes to discard certain information.

Additionally, the task of guessing a piece of information (equivalent to one outcome of $\mathcal{X}$) is no less challenging than guessing a bit-string of length $H(\mathcal{X})$, disregarding momentarily the effort involved in decoding the guess.

\subsection{Minimum entropy}
\begin{definition}[\textit{Min-entropy}]
    The min-entropy is a measure of the most conservative assessment of the unpredictability of a set of outcomes.
    It is defined for $\mathcal{X}$ as:
    \[H_\infty(\mathcal{X})=-\log(\max_i\Pr_i)\]
\end{definition}
In essence, it represents the entropy of a random variable with a uniform distribution, where each outcome has a probability of $\max_ip_i$.

It's worth noting that guessing the most common outcome of $\mathcal{X}$ is no less challenging than guessing a bit-string of length $H_\infty(\mathcal{X})$.
\begin{example}
    Consider the random variable $\mathcal{X}$ defined as:
    \[\mathcal{X}=\begin{cases}
        0^{128}  \qquad \qquad          \text{with probability} \frac{1}{2}\\
        a        \quad \qquad \qquad    \text{with probability} \frac{1}{2^{128}}
    \end{cases}\]
    Here, $a \in 1\{0,1\}^{127}$.
    Predicting an outcome shouldn't be too difficult: just predict $0^{128}$:
    \[H(\mathcal{X})=\dfrac{1}{2}\left(-\log_2\left(\dfrac{1}{2}\right)\right)+2^{127}\frac{1}{2^{128}}\left(-\log_2\left(\frac{1}{2^{128}}\right)\right)= 64.5b\]
    \[H_\infty(\mathcal{X})=-\log_2\left(\dfrac{1}{2}\right)=1b\]
    Min-entropy indicates that guessing the most common output is as difficult as guessing a single bit string.
\end{example}